% ─────────────────────────────────────────────────────────────────────
%  A Layered, Axiom-Isolated Lean 4 Formalization of the Congruent
%  Number Problem up to Tunnell's Criterion
%
%  Manuscript for internal peer review — April 2026
%  Repository: congruent-relay, commit 5b42d4b
% ─────────────────────────────────────────────────────────────────────
\documentclass[11pt,a4paper]{article}

% ── Packages ──────────────────────────────────────────────────────────
\usepackage[utf8]{inputenc}
\usepackage[T1]{fontenc}
\usepackage{lmodern}
\usepackage{amsmath,amssymb,amsthm}
\usepackage{mathtools}
\usepackage{booktabs}
\usepackage{enumitem}
\usepackage{hyperref}
\usepackage[margin=1in]{geometry}
\usepackage{xcolor}
\usepackage{listings}
\usepackage{microtype}
\usepackage{cleveref}

% ── Theorem environments ──────────────────────────────────────────────
\newtheorem{theorem}{Theorem}[section]
\newtheorem{lemma}[theorem]{Lemma}
\newtheorem{corollary}[theorem]{Corollary}
\newtheorem{proposition}[theorem]{Proposition}
\theoremstyle{definition}
\newtheorem{definition}[theorem]{Definition}
\newtheorem{remark}[theorem]{Remark}
\newtheorem{axiomenv}[theorem]{Axiom}

% ── Listings (Lean 4) ────────────────────────────────────────────────
\lstdefinelanguage{Lean4}{
  morekeywords={def,theorem,lemma,axiom,opaque,namespace,end,where,
    let,if,then,else,match,with,by,exact,simp,intro,obtain,have,
    apply,constructor,rfl,omega,nlinarith,linarith,ext,fun,
    Finset,Int,Nat,Prop,Set,Even,True,sorry,induction},
  sensitive=true,
  morecomment=[l]{--},
  morecomment=[n]{/-}{-/},
  morestring=[b]",
}
\lstset{
  language=Lean4,
  basicstyle=\ttfamily\small,
  keywordstyle=\bfseries\color{blue!70!black},
  commentstyle=\itshape\color{gray},
  breaklines=true,
  frame=single,
  framerule=0.4pt,
  xleftmargin=1em,
  aboveskip=0.8em,
  belowskip=0.5em,
  extendedchars=false,
}

% ── Macros ────────────────────────────────────────────────────────────
\newcommand{\Z}{\mathbb{Z}}
\newcommand{\Q}{\mathbb{Q}}
\newcommand{\N}{\mathbb{N}}
\newcommand{\lean}[1]{\texttt{#1}}
\newcommand{\solSet}{\mathrm{solSet}}
\newcommand{\repCount}{\mathrm{repCount}}
\newcommand{\solSetInf}{\mathrm{solSet}^{\infty}}
\newcommand{\BSDfor}{\mathrm{BSD\_for}}
\newcommand{\En}{E_n}

% ─────────────────────────────────────────────────────────────────────
\title{%
  A Layered, Axiom-Isolated Lean~4 Formalization of the \\
  Congruent Number Problem up to Tunnell's Criterion%
}
\author{Papanokechi}
\date{April 2026 \quad---\quad Internal Review Draft}

\begin{document}
\maketitle

% ═════════════════════════════════════════════════════════════════════
\begin{abstract}
We present a 954-line Lean~4 formalization of the combinatorial and
structural backbone of Tunnell's criterion for the congruent number
problem.  The development is organized into six layers, each
mathematically natural and fully proved, culminating in a clean axiom
boundary where the sole remaining assumptions are the Birch and
Swinnerton-Dyer (BSD) conjecture and Tunnell's conditional theorem.

The formalization contains zero \lean{sorry} placeholders.  The only
axiom is \lean{tunnell\_conditional\_on\_BSD}, which encodes
Tunnell's 1983 result.  The BSD hypothesis is represented as
$\BSDfor(n) \coloneqq \mathrm{AlgebraicRank}(n) =
\mathrm{AnalyticRank}(n)$, where both rank functions are opaque
definitions, preventing vacuous discharge.

Key contributions include a generic fixed-point-free involution lemma
on \lean{Finset} (proved by strong induction), an eightfold sign
symmetry for positive-definite ternary quadratic forms, a bridging
theorem connecting bounded and unbounded solution sets, and a
theta-series interface that recovers Tunnell's counting quantities.

The formal development is paired with a Python computational pipeline
(98 tests, validated against OEIS~A003273) that verifies Tunnell's
criterion for all squarefree~$n \le 500$.
\end{abstract}

% ═════════════════════════════════════════════════════════════════════
\section{Introduction}\label{sec:intro}

A positive integer $n$ is a \emph{congruent number} if it is the area
of a right triangle with rational side lengths.  Equivalently, $n$ is
congruent if and only if the elliptic curve
\[
  \En : y^2 = x^3 - n^2 x
\]
has positive Mordell--Weil rank over~$\Q$.  The congruent number
problem---determining which integers are congruent---is one of the
oldest unsolved problems in number theory, with roots in Arab
manuscripts from the tenth century~\cite{Koblitz1993}.

Tunnell's theorem~\cite{Tunnell1983} provides an effective criterion:
a squarefree positive integer $n$ is congruent if and only if certain
representation counts of ternary quadratic forms satisfy an equality
condition, \emph{provided} the Birch and Swinnerton-Dyer conjecture
holds for~$\En$.  The forward direction is unconditional; the reverse
direction requires BSD.

\paragraph{Why formalize?}
Tunnell's criterion involves a delicate interplay between
combinatorial counting (ternary quadratic forms), analytic number
theory (half-integral weight modular forms), and deep conjectures
(BSD\@).  A formalization makes the dependency structure explicit and
machine-checkable.  However, formalizing the full analytic machinery
is currently out of reach for any proof assistant.  Our strategy is
to formalize everything that \emph{can} be proved combinatorially and
isolate the analytic content behind a single, honest axiom.

\paragraph{Design philosophy.}
\begin{enumerate}[nosep]
  \item \textbf{Layered architecture.}  Six sections, each building
    on the previous, with a clear separation between combinatorial
    content (proved) and analytic content (axiomatized).
  \item \textbf{Axiom isolation.}  The BSD conjecture enters as a
    single, mathematically meaningful definition---not a vacuous
    \lean{True} placeholder.
  \item \textbf{Reproducibility.}  All definitions mirror a Python
    implementation; a test suite validates numerical agreement.
\end{enumerate}

\paragraph{Contributions.}
\begin{itemize}[nosep]
  \item A complete Lean~4 formalization of bound independence,
    sign symmetry, and the unbounded bridging lemma for
    positive-definite ternary quadratic forms.
  \item A generic \lean{Finset} involution lemma
    (\lean{card\_even\_of\_involution}) proved by strong induction,
    yielding even cardinality of nonzero solution sets.
  \item A theta-series bridge connecting combinatorial representation
    counts to Tunnell's criterion quantities.
  \item A clean axiom boundary: one axiom, two opaque definitions,
    zero \lean{sorry}.
  \item A paired Python pipeline (98~tests) validating
    Tunnell's criterion for $n \le 500$.
\end{itemize}

% ═════════════════════════════════════════════════════════════════════
\section{Related Work}\label{sec:related}

Formalizations of elliptic curves in proof assistants remain rare.
Bartzia and Strub~\cite{BartziaStrub2014} formalized elliptic curve
arithmetic in Coq with a focus on cryptographic applications.
The Mathlib library for Lean~4~\cite{Mathlib2024} includes definitions
of elliptic curves and Weierstrass models, but does not yet cover
modular forms of half-integral weight or the Shimura correspondence
that underlies Tunnell's proof.

The congruent number problem has been treated computationally by many
authors.  Brent and van~de~Lune~\cite{BrentVanDeLune2007} verified
congruent numbers up to~$10^6$ using a combination of analytic and
algebraic methods.  Our contribution is orthogonal: we do not push
computational boundaries but instead provide a machine-checked proof
that the combinatorial layer of Tunnell's criterion is correctly
implemented.

The strategy of isolating deep conjectures as explicit axioms has
precedents.  Buzzard et al.~\cite{BuzzardEtAl2020} advocate for
formalizations that clearly delineate the ``sorry boundary'' and
argue that even partial formalizations provide value by exposing
implicit assumptions.  Our work follows this philosophy: the BSD
conjecture is axiomatized as a non-trivial proposition
$\BSDfor(n)$, making the dependency visible and auditable.

% ═════════════════════════════════════════════════════════════════════
\section{Overview of the Lean Development}\label{sec:lean}

The formalization resides in a single file,
\texttt{formal/CongruentStubs.lean} (954~lines), within the
\lean{Congruent} namespace.  It depends on Mathlib~4 for
\lean{Finset}, \lean{Set}, \lean{Int}, \lean{Rat}, and standard
algebraic infrastructure.  We now describe each of the six sections.

% ─────────────────────────────────────────────────────────────────────
\subsection{Preliminaries: Quadratic Forms and Solution Sets}\label{sec:s1}

\begin{definition}[Solution set]\label{def:solset}
Let $a, b, c, t \in \Z$ and $B \in \N$.  The \emph{bounded solution
set} is
\[
  \solSet(a,b,c,t,B) \;=\;
  \bigl\{(x,y,z) \in \Z^3 : |x|,|y|,|z| \le B
    \text{ and } ax^2 + by^2 + cz^2 = t\bigr\}.
\]
\end{definition}

This is defined as a \lean{Finset (Int \texttimes{} Int \texttimes{} Int)}
via filtering the product of integer ranges.

\begin{definition}[Representation count]\label{def:repcount}
\[
  \repCount(a,b,c,t) \;=\;
  \begin{cases}
    0 & \text{if } t < 0,\\
    \bigl|\solSet(a,b,c,t,|t|)\bigr| & \text{otherwise.}
  \end{cases}
\]
\end{definition}

The bound $B = |t|$ is generous but correct for positive-definite
forms; \S\ref{sec:s2} proves that any $B \ge |t|$ yields the same
count.

\begin{definition}[Squarefree]\label{def:sqfree}
A natural number $n$ is \emph{squarefree} if $n > 0$ and no
$d^2$ with $d > 1$ divides~$n$.
\end{definition}

% ─────────────────────────────────────────────────────────────────────
\subsection{Bound Independence and Stability}\label{sec:s2}

For positive-definite forms ($a, b, c > 0$) with $t \ge 0$, any
solution $(x,y,z)$ satisfies $|x| \le |t|$ (and similarly for $y$
and $z$), because $ax^2 \le ax^2 + by^2 + cz^2 = t$.  This gives:

\begin{theorem}[Coordinate bound]\label{thm:coord-bound}
If $a > 0$ and $ax^2 \le t$ with $t \ge 0$, then $|x| \le |t|$.
\end{theorem}

\begin{theorem}[Monotonicity]\label{thm:mono}
If $B_1 \le B_2$ then
$\solSet(a,b,c,t,B_1) \subseteq \solSet(a,b,c,t,B_2)$.
\end{theorem}

\begin{theorem}[Stability]\label{thm:stable}
For $a, b, c > 0$ and $t \ge 0$, if $B \ge |t|$ then
$\solSet(a,b,c,t,B) = \solSet(a,b,c,t,|t|)$.  In particular,
$\repCount$ is independent of the choice of sufficient bound.
\end{theorem}

% ─────────────────────────────────────────────────────────────────────
\subsection{Sign Symmetry and the Unbounded Bridge}\label{sec:s3}

\begin{theorem}[Eightfold sign symmetry]\label{thm:sign-sym}
For $a, b, c > 0$, $t \ge 0$, and $B \ge |t|$, if
$(x,y,z) \in \solSet(a,b,c,t,B)$ then all eight sign flips
$(\pm x, \pm y, \pm z) \in \solSet(a,b,c,t,B)$.
\end{theorem}

\begin{theorem}[Origin characterization]\label{thm:origin}
$(0,0,0) \in \solSet(a,b,c,t,B)$ if and only if $t = 0$ (assuming $B > 0$).
\end{theorem}

\begin{definition}[Unbounded solution set]
\[
  \solSetInf(a,b,c,t) \;=\;
  \bigl\{(x,y,z) \in \Z^3 : ax^2 + by^2 + cz^2 = t\bigr\}.
\]
\end{definition}

\begin{theorem}[Finiteness]\label{thm:finite}
For positive-definite forms with $t \ge 0$,
$\solSetInf(a,b,c,t)$ is finite.
\end{theorem}

\begin{theorem}[Bridging lemma]\label{thm:bridge}
For $a, b, c > 0$ and $t \ge 0$,
\[
  \solSet(a,b,c,t,|t|) \;=\; \solSetInf(a,b,c,t)
\]
as subsets of $\Z^3$.  Hence $\repCount(a,b,c,t) =
|\solSetInf(a,b,c,t)|$.
\end{theorem}

This theorem is the key correctness property: the computable
\lean{repCount} counts the same solutions as the mathematical
definition.

% ─────────────────────────────────────────────────────────────────────
\subsection{Orbit Pairing and Parity}\label{sec:s4}

The central technical result is a generic involution lemma on finite
sets:

\begin{lemma}[Fixed-point-free involution $\Rightarrow$ even cardinality]
\label{lem:involution}
Let $S$ be a finite set and $f\colon S \to S$ an involution
($f \circ f = \mathrm{id}$) with no fixed points.  Then
$|S|$ is even.
\end{lemma}

\begin{proof}[Proof sketch]
By strong induction on $|S|$.  If $S = \varnothing$ then
$|S| = 0$ is even.  Otherwise, pick $x \in S$.  Since $f$ has no
fixed points, $f(x) \ne x$ and both $x, f(x) \in S$.  Let
$S' = S \setminus \{x, f(x)\}$.  The involution restricts to $S'$
(if $y \in S'$ then $f(y) \ne x$ and $f(y) \ne f(x)$, so
$f(y) \in S'$).  By the inductive hypothesis, $|S'|$ is even, so
$|S| = |S'| + 2$ is even.
\end{proof}

This is formalized as \lean{card\_even\_of\_involution} and
proved entirely within Lean~4 using \lean{Finset.strongInduction}.

\begin{definition}[Nonzero solution set]
$\mathrm{solSetNZ}(a,b,c,t,B) = \solSet(a,b,c,t,B) \setminus \{(0,0,0)\}$.
\end{definition}

\begin{corollary}\label{cor:even-nonzero}
The negation map $(x,y,z) \mapsto (-x,-y,-z)$ is a
fixed-point-free involution on $\mathrm{solSetNZ}$.  Hence
$|\mathrm{solSetNZ}|$ is even.
\end{corollary}

\begin{theorem}[Parity decomposition]\label{thm:parity}
\[
  \repCount(a,b,c,t) \;=\;
  |\mathrm{solSetNZ}| + \begin{cases} 1 & \text{if } t = 0,\\ 0 & \text{if } t > 0.\end{cases}
\]
In particular, $\repCount(a,b,c,t)$ is even when $t > 0$.
\end{theorem}

% ─────────────────────────────────────────────────────────────────────
\subsection{Theta-Series Bridge}\label{sec:s5}

We introduce the combinatorial alias:

\begin{definition}[Theta coefficient]
$\mathrm{thetaCoeff}(a,b,c,t) \coloneqq \repCount(a,b,c,t)$.
\end{definition}

This naming reflects the classical interpretation: the $n$-th
coefficient of the theta series
$\theta_{a,b,c}(q) = \sum_{x,y,z \in \Z} q^{ax^2+by^2+cz^2}$
is exactly $\repCount(a,b,c,n)$.

\begin{theorem}[Theta--solution set agreement]\label{thm:theta-sol}
$\mathrm{thetaCoeff}(a,b,c,t) = |\solSetInf(a,b,c,t)|$ for
positive-definite forms with $t \ge 0$.
\end{theorem}

\begin{definition}[Tunnell counts]\label{def:tunnell-counts}
For an integer $n$:
\[
  \mathrm{tunnellCounts}(n) =
  \begin{cases}
    \bigl(\repCount(2,1,8,n),\; \repCount(2,1,32,n)\bigr)
      & \text{if } n \text{ is odd},\\[4pt]
    \bigl(\repCount(4,1,8,n/2),\; \repCount(4,1,32,n/2)\bigr)
      & \text{if } n \text{ is even}.
  \end{cases}
\]
\end{definition}

\begin{definition}[Tunnell predicate]\label{def:tunnell-pred}
$\mathrm{isCongruentTunnell}(n)$ holds iff $A = 2B$ where
$(A, B) = \mathrm{tunnellCounts}(n)$.
\end{definition}

\begin{theorem}\label{thm:theta-tunnell}
The theta-series formulation and the direct counting formulation
coincide: $\mathrm{tunnellThetaCoeffs}(n) = \mathrm{tunnellCounts}(n)$.
\end{theorem}

% ─────────────────────────────────────────────────────────────────────
\subsection{Tunnell's Criterion Under BSD}\label{sec:s6}

\begin{definition}[Congruent number]
A natural number $n$ is \emph{congruent} if $n > 0$ and there exist
positive rationals $a, b, c$ with $a^2 + b^2 = c^2$ and
$ab/2 = n$.
\end{definition}

\begin{definition}[Elliptic curve $\En$]
For a natural number $n$, the predicate $\mathrm{OnCurveEn}(n,x,y)$
holds iff $y^2 = x^3 - n^2 x$ (over~$\Q$).
\end{definition}

To give the BSD hypothesis a mathematically meaningful type, we
introduce opaque rank functions:

\begin{definition}[Rank functions (opaque)]
\begin{align*}
  &\mathrm{AlgebraicRank}(n) : \N && \text{(Mordell--Weil rank of $\En(\Q)$)}\\
  &\mathrm{AnalyticRank}(n) : \N && \text{(order of vanishing of $L(\En,s)$ at $s=1$)}
\end{align*}
These are declared \lean{opaque} in Lean~4: their definitions are
irreducible, so no properties can be derived without the axiom below.
\end{definition}

\begin{definition}[BSD hypothesis]
\[
  \BSDfor(n) \;\coloneqq\;
  \mathrm{AlgebraicRank}(n) = \mathrm{AnalyticRank}(n).
\]
\end{definition}

\begin{remark}
This is a genuine mathematical proposition---it cannot be discharged
by \lean{trivial} or \lean{exact trivial}.  A reviewer can verify
that the reverse direction of Tunnell's theorem requires a
non-vacuous hypothesis.
\end{remark}

The only axiom in the entire formalization is:

\begin{axiomenv}[Tunnell's theorem, conditional on BSD]\label{ax:tunnell}
For every squarefree $n \in \N$:
\begin{enumerate}[nosep]
  \item \textbf{Forward (unconditional):}
    $\mathrm{IsCongruentNumber}(n) \Rightarrow
     \mathrm{isCongruentTunnell}(n)$.
  \item \textbf{Reverse (conditional):}
    $\BSDfor(n) \Rightarrow
     \mathrm{isCongruentTunnell}(n) \Rightarrow
     \mathrm{IsCongruentNumber}(n)$.
\end{enumerate}
\end{axiomenv}

\begin{corollary}[Forward direction]\label{cor:forward}
If $n$ is congruent and squarefree, the Tunnell count condition holds.
\end{corollary}

\begin{corollary}[Reverse direction under BSD]\label{cor:reverse}
If $\BSDfor(n)$ holds and the Tunnell count condition is satisfied
for squarefree~$n$, then $n$ is congruent.
\end{corollary}

Both corollaries are one-line consequences of the axiom, formalized
as \lean{tunnell\_forward} and \lean{tunnell\_reverse\_under\_BSD}.

% ═════════════════════════════════════════════════════════════════════
\section{Axiom Boundary}\label{sec:boundary}

The design principle is \emph{minimal axiomatization}: everything
that can be proved combinatorially is proved; the single remaining
axiom encodes a deep theorem whose proof requires analytic machinery
beyond the reach of current proof assistants.

\begin{table}[ht]
\centering
\begin{tabular}{@{}lll@{}}
\toprule
\textbf{Component} & \textbf{Status} & \textbf{Lean construct} \\
\midrule
Combinatorial lemmas (\S\S\ref{sec:s1}--\ref{sec:s4})
  & Fully proved & \lean{theorem} \\
Involution lemma
  & Fully proved & \lean{theorem} \\
Theta-series bridge
  & Fully proved & \lean{theorem} \\
BSD hypothesis
  & \textbf{Single axiom} & \lean{def BSD\_for} \\
Rank functions
  & \textbf{Opaque defs} & \lean{opaque} \\
Tunnell's theorem
  & \textbf{Single axiom} & \lean{axiom} \\
\bottomrule
\end{tabular}
\caption{Axiom boundary.  Everything above the line is fully
machine-checked; everything below is explicitly axiomatized.}
\label{tab:boundary}
\end{table}

\paragraph{Why \lean{opaque} and not \lean{axiom}?}
The rank functions \lean{AlgebraicRank} and \lean{AnalyticRank} are
mathematical quantities that \emph{exist} unconditionally---the
Mordell--Weil rank is a well-defined natural number.  They are
\lean{opaque} (not \lean{axiom}) because we do not assert any
property about them; we merely need them as types.  The BSD
conjecture is then a \lean{def} producing a \lean{Prop}, and
Tunnell's theorem is the sole \lean{axiom}.

This three-level separation (\lean{opaque} quantities $\to$
\lean{def} hypothesis $\to$ \lean{axiom} theorem) makes the
dependency structure maximally transparent.

% ═════════════════════════════════════════════════════════════════════
\section{Evaluation}\label{sec:eval}

\begin{table}[ht]
\centering
\begin{tabular}{@{}lr@{}}
\toprule
\textbf{Metric} & \textbf{Value} \\
\midrule
Lines of Lean~4            & 954 \\
\lean{sorry} count         & 0 \\
\lean{axiom} count         & 1 \\
\lean{opaque} count        & 2 \\
Theorems / lemmas          & ${\sim}30$ \\
Definitions                & ${\sim}15$ \\
Python test suite          & 98 pass, 1 skip \\
Computational range        & $n \le 500$ (squarefree) \\
Congruent numbers found    & 185 / 306 \\
\bottomrule
\end{tabular}
\caption{Project metrics at commit \texttt{5b42d4b}.}
\label{tab:metrics}
\end{table}

\paragraph{Reproducibility.}
The Lean formalization builds with Mathlib~4.  The Python pipeline
runs as:
\begin{lstlisting}[language=bash,keywordstyle={},commentstyle={}]
pytest -q                       # 98 pass, 1 skip
python -m src.runner --max-n 500  # full scan
\end{lstlisting}
A \texttt{Dockerfile} is provided for containerized reproduction.

\paragraph{Proof techniques.}
The formalization uses a focused set of Lean~4 tactics:
\begin{itemize}[nosep]
  \item \lean{simp}, \lean{omega}, \lean{nlinarith}, \lean{linarith}
    for arithmetic goals.
  \item \lean{ext} for set/function extensionality.
  \item \lean{Finset.strongInduction} for the involution lemma.
  \item No \lean{decide} on large instances; no \lean{native\_decide}.
\end{itemize}

% ═════════════════════════════════════════════════════════════════════
\section{Conclusion and Future Work}\label{sec:conclusion}

We have presented a layered, axiom-isolated Lean~4 formalization of
the combinatorial core of Tunnell's criterion for the congruent
number problem.  The development achieves zero \lean{sorry}
placeholders with a single, honest axiom encoding Tunnell's
conditional theorem under BSD\@.

\paragraph{Limitations.}
The formalization does not cover modular forms of half-integral
weight, the Shimura correspondence, or the analytic continuation
of $L$-functions---all of which are needed for Tunnell's
\emph{proof}.  We formalize Tunnell's \emph{statement}, not its
proof.

\paragraph{Future directions.}
\begin{enumerate}[nosep]
  \item \textbf{Modular forms.}  Formalize weight-$3/2$ modular
    forms and the Shimura lift, which would allow replacing the
    axiom with a proved theorem (modulo BSD).
  \item \textbf{Analytic rank.}  Connect $\mathrm{AnalyticRank}$
    to a formalized $L$-function, enabling partial BSD results
    (e.g., Coates--Wiles for CM curves).
  \item \textbf{Heegner points.}  Formalize the Gross--Zagier
    theorem to give computational access to rank~1 cases.
  \item \textbf{Mathlib integration.}  Contribute the involution
    lemma and quadratic-form infrastructure to Mathlib.
  \item \textbf{Extended computation.}  Push the Python pipeline
    to $n \le 10^6$ with optimized enumeration.
\end{enumerate}

% ═════════════════════════════════════════════════════════════════════
\appendix
\section{File Tour}\label{app:tour}

Table~\ref{tab:tour} summarizes the structure of
\texttt{formal/CongruentStubs.lean}.

\begin{table}[ht]
\centering
\small
\begin{tabular}{@{}llp{6cm}@{}}
\toprule
\textbf{Section} & \textbf{Lines} & \textbf{Key items} \\
\midrule
\S1 Foundations
  & $\sim$66--107
  & \lean{IsSquarefree}, \lean{solSet}, \lean{repCount} \\
\S2 Basic properties
  & $\sim$107--182
  & \lean{repCount\_neg}, \lean{repCount\_nonneg},
    \lean{solSet\_neg\_x/y/z} \\
\S3a Bound independence
  & $\sim$182--287
  & \lean{solSet\_mono}, \lean{solSet\_stable},
    \lean{repCount\_stable} \\
\S3b Sign symmetry
  & $\sim$287--341
  & \lean{solSet\_sign\_symmetry}, \lean{solSet\_origin} \\
\S3c Unbounded bridge
  & $\sim$341--414
  & \lean{solSet\textsuperscript{$\infty$}},
    \lean{solSet\textsuperscript{$\infty$}\_finite},
    \lean{solSet\_eq\_solSet\textsuperscript{$\infty$}} \\
\S4 Orbit pairing
  & $\sim$414--690
  & \lean{card\_even\_of\_involution},
    \lean{repCount\_even\_of\_pos} \\
Tunnell defs
  & $\sim$690--754
  & \lean{tunnellCounts}, \lean{isCongruentTunnell},
    \lean{tunnellCounts\_cases} \\
\S5 Theta series
  & $\sim$754--849
  & \lean{thetaCoeff},
    \lean{thetaCoeff\_eq\_solSet\textsuperscript{$\infty$}\_card} \\
\S6 Tunnell axiom
  & $\sim$849--954
  & \lean{IsCongruentNumber}, \lean{BSD\_for},
    \lean{tunnell\_conditional\_on\_BSD} \\
\bottomrule
\end{tabular}
\caption{File tour for \texttt{CongruentStubs.lean} (954~lines).}
\label{tab:tour}
\end{table}

% ─────────────────────────────────────────────────────────────────────
\section{Selected Lean Excerpts}\label{app:code}

\paragraph{Core definition: solution set.}
\begin{lstlisting}
def solSet (a b c t : Int) (B : Nat) :
    Finset (Int * Int * Int) :=
  let range := Finset.Icc (-(B : Int)) (B : Int)
  (range.product (range.product range)).filter
    fun (x, y, z) =>
      a * x * x + b * y * y + c * z * z = t
\end{lstlisting}

\paragraph{Involution lemma (statement).}
\begin{lstlisting}
theorem card_even_of_involution [DecidableEq a]
    (s : Finset a) (f : a -> a)
    (hf_mem : forall x in s, f x in s)
    (hf_inv : forall x in s, f (f x) = x)
    (hf_fix : forall x in s, f x /= x) :
    Even s.card
\end{lstlisting}

\paragraph{BSD hypothesis and Tunnell axiom.}
\begin{lstlisting}
opaque AlgebraicRank (n : Nat) : Nat
opaque AnalyticRank (n : Nat) : Nat

def BSD_for (n : Nat) : Prop :=
  AlgebraicRank n = AnalyticRank n

axiom tunnell_conditional_on_BSD :
  forall n : Nat, IsSquarefree n ->
    (IsCongruentNumber n -> isCongruentTunnell n) /\\
    (BSD_for n -> isCongruentTunnell n
                -> IsCongruentNumber n)
\end{lstlisting}

\paragraph{Bridging lemma.}
\begin{lstlisting}
theorem solSet_eq_solSetInf {a b c t : Int}
    (ha : a > 0) (hb : b > 0) (hc : c > 0)
    (ht : t >= 0) :
    (solSet a b c t t.natAbs : Set _)
      = solSetInf a b c t
\end{lstlisting}

% ═════════════════════════════════════════════════════════════════════
\bibliographystyle{plain}
\begin{thebibliography}{99}

\bibitem{Tunnell1983}
J.\,B.~Tunnell,
``A classical Diophantine problem and modular forms of weight~3/2,''
\emph{Inventiones mathematicae}, vol.~72, pp.~323--334, 1983.

\bibitem{Koblitz1993}
N.~Koblitz,
\emph{Introduction to Elliptic Curves and Modular Forms},
2nd~ed., Springer, 1993.

\bibitem{BartziaStrub2014}
E.~Bartzia and P.-Y.~Strub,
``A formal library for elliptic curves in the Coq proof assistant,''
in \emph{Proc.\ ITP}, 2014, pp.~77--92.

\bibitem{Mathlib2024}
The Mathlib Community,
``Mathlib4: The Lean 4 Mathematical Library,''
\url{https://github.com/leanprover-community/mathlib4}, 2024.

\bibitem{BrentVanDeLune2007}
R.\,P.~Brent and J.~van de~Lune,
``A note on the computation of congruent numbers,''
manuscript, 2007.

\bibitem{BuzzardEtAl2020}
K.~Buzzard, J.~Commelin, and P.~Massot,
``Formalising perfectoid spaces,''
in \emph{Proc.\ CPP}, 2020, pp.~299--312.

\end{thebibliography}

\end{document}
